\documentclass[11pt]{article}
\usepackage{amsmath,amssymb}
\title{Why the original endpoint minor needs contact cancellation before longest-time decay}\author{}\date{}
\begin{document}\maketitle
\paragraph{Input and target.}
In the actual cubic curl and endpoint minor calculation, the exact plaquette cubic grouping, factor-24 all-cross
contraction, endpoint minor and transported curvature vertex retain their
stated partners. This note tests the proposed extra $1/(1+T)$ on that
minor itself. It proves that such a bound cannot hold uniformly for an
individual charged scalar minor, even deep inside the ORIGINAL finite
relative grid. The short-line contact must be canceled with the complete
quartic/ghost/Haar and source partners first. This is not a divergence
claim for their full sum or for the original field theory.

\paragraph{1. The exact short-line term.}
Polarize the endpoint determinant of the actual cubic curl and endpoint minor calculation. For two heat times $s,T$
and a neighboring endpoint pair $x,x+\hat\mu$, let
\[
 \mathcal B(P_s,P_T)=\tfrac12[
 (P_s)_{00}(P_T)_{11}+(P_T)_{00}(P_s)_{11}
 -(P_s)_{01}(P_T)_{10}-(P_T)_{01}(P_s)_{10}].                       \tag{1}
\]
Integrating (1) over both positive times is exactly the Green endpoint
minor, not a new vertex.
Choose an odd cube side $R=2m+1$, a $\nu$ component with $\nu\ne\mu$,
and the central pair with $\mu$ positions $m,m+1$. The original product
grid is invariant under their reflection. At zero field the two endpoint
matrices are therefore
\[
 P_s=\begin{pmatrix}u&v\\v&u\end{pmatrix},\qquad
 P_T=\begin{pmatrix}U&V\\V&U\end{pmatrix}.
\]
Positivity gives $v\ge0$ and $U-V\ge0$. Hence
\[
 \boxed{\mathcal B=(u-v)U+v(U-V)\ge(u-v)U.}                       \tag{2}
\]
For a sufficiently small fixed $s>0$, the no-jump contribution gives
$u\ge e^{-8s}$, while $v\le1-e^{-8s}\le8s$. Thus $u-v\ge c>0$
uniformly in the original cube size. Central product eigenfunctions give
$U\ge cT^{-2}$ for $1\le T\le cR^2$; this also follows by retaining the
low-frequency sine/cosine modes with indices at most $cR/\sqrt T$.
Therefore the minor can be of order $T^{-2}$, not $T^{-3}$.
The two literal differences in the vertex can both act on the SHORT
line. Their algebraic presence does not pay a longest-time derivative.
The exact dangerous piece in (2) is the short-link Dirichlet energy
$u-v$ multiplying the undifferentiated long heat diagonal $U$.

\paragraph{2. The geometric current has the same short-time contact.}
Put the actual magnetic scalar connection of strength $b=2/R^2$ into (1).
At $s=0$ its scalar current part is exactly
\[
 \mathcal B_\epsilon(P_0,P_T)
       =\tfrac12[p_T^\epsilon(x,x)+p_T^\epsilon(x+\hat\mu,x+\hat\mu)].
                                                                    \tag{3}
\]
Thus minus its second flux derivative is the average of the two original
closed-loop area second moments. The derivative is NOT zero merely
because the endpoint minor contains differences.

Here is a self-contained quantitative lower bound, avoiding an assumption
about reflected bridge areas. First consider the unrestricted rate-one
nearest-neighbor walk on $\mathbb Z^4$. Conditional on return at time $T$,
the numbers $n_j$ of positive and negative jumps in each coordinate agree;
the coordinates' counts are independent, and one count has probabilities
proportional to $T^{2n}/(n!)^2$. Shifting this series gives
\[
 E n^2=T^2,\qquad E n^4=T^2(T^2+2En+1)\le4T^4\quad(T\ge1).
\]
The elementary second-moment inequality for $n^2$ then implies
$P(n\ge T/\sqrt2)\ge1/16$, hence $En\ge cT$.

Conditional on all counts, the order of the jumps is a uniform
permutation. For its projected closed 01 polygon let $A$ be the signed
lattice area. If its labeled steps are $v_i$ and
$c_{ij}=v_i\wedge v_j$, then
$A=\frac12\sum_{i<j}c_{ij}\operatorname{sign}(U_j-U_i)$ for independent
continuous ordering priorities $U_i$. Disjoint sign pairs have mean zero;
shared-endpoint pairs have correlation $\pm1/3$.
Since $\sum_jc_{ij}=0$ for a closed word, expansion gives
\[
 E[A^2\mid(n_j)]=\frac1{12}\sum_{i<j}c_{ij}^2
                         =\frac{n_0n_1}{3}.                     \tag{4}
\]
Other coordinate steps do not change the relative order distribution.
It follows that the unrestricted return-weighted current satisfies
\[
 E[(bA)^2;X_T=x]\ge cb^2 T^2 p_T^{\mathbb Z^4}(0)\ge cb^2.
                                                                    \tag{5}
\]
The last return bound is elementary: the one-dimensional Fourier integral
of $e^{-2T(1-\cos k)}$ is bounded below by $cT^{-1/2}$ by integrating
$|k|\le T^{-1/2}$, and the four coordinates factorize.

Transfer this lower bound to the ORIGINAL finite grid at a central
endpoint by choosing $T=R^{1/3}$. Every path with fewer than $R/4$ total
jumps stays away from its boundary, so its original killed/reflected
weight and magnetic current equal the unrestricted ones. The total
number $N$ of jumps is Poisson with mean $8T$, and $|A|^2\le CN^4$.
The discarded area moment is at most
\[
 Cb^2E[N^4;N\ge R/4]
 \le Cb^2(1+T)^4\exp[-R/4+CT]=o(b^2).                            \tag{6}
\]
For example (6) follows by inserting $e^{N-R/4}$ and differentiating the
Poisson exponential generating function four times. Thus (5), with a
smaller constant, holds at the two central ORIGINAL-grid endpoints.
Combining with (3) proves
\[
 -\partial_\epsilon^2\mathcal B_\epsilon(P_0,P_T)\big|_0
       \ge cb^2,\qquad T=R^{1/3}\longrightarrow\infty.            \tag{7}
\]
A diagramwise attempt to attach an extra $1/(1+T)$ to the natural
current scale $b^2T^2(1+T)^{-2}$ would instead give $Cb^2/(1+T)$,
contradicting (7). Continuity in $s$ rules out such a uniform bound on
all positive short times as well. Equation(7) concerns the geometric
scalar-current part; actual differentiated vertex, spin and ghost terms
may cancel it and have NOT been deleted from the full coefficient.

\paragraph{3. Exact checks and the next estimate.}
The product-spectrum finite calculation evaluates (2) on original cubes
$R=9,17,33,65,129$, at $s=1/64$ and $T=R^{1/3}$.
The ratio $\mathcal B/(uU)$ stays above $0.985$ while
$T\mathcal B/(uU)$ grows. Exact enumeration of closed words with
$(n_0,n_1)=(1,1),(1,2),(2,2)$ gives area variances
$1/3,2/3,4/3$, checking (4).

The next valid main estimate is on the FULL contact-subtracted,
Ward-completed current. It must match the local short-line term in (2)
and the zero-time current contact (3) with the original quartic/Haar/
ghost and endpoint/spin partners. The derivative network contact and boundary defect calculation exhibits one scalar
contraction where this cancellation is exact, with a retained mixed
boundary defect. The finite boundary nonlinear ward cancellation calculation supplies the complete boundary mass Ward.
Neither identity yet proves that every tensor sector of the actual cubic curl and endpoint minor calculation has
been grouped correctly. A uniform bound on the uncompleted minor is
not a substitute for that missing matching step.
No obstruction to a renormalized construction or S94 is inferred from
this necessary correction to the attempted diagramwise bound.
\end{document}
